%====================== REFERENCIAL TEÓRICO ======================
\section{Referencial teórico}

\subsection{A Língua Brasileira de Sinais (Libras)}

A Libras é uma língua natural, de modalidade visuoespacial, dotada de estrutura
gramatical própria — e não uma codificação gestual do português, equívoco comum
entre pessoas não familiarizadas com a língua de sinais (QUADROS; KARNOPP, 2004;
QUADROS, 2019). Segundo Quadros e Karnopp (2004), a estrutura das línguas de
sinais, a partir dos estudos pioneiros de Stokoe, é composta por cinco parâmetros
fonológicos que, isoladamente, não carregam significado, mas que em combinação
formam os sinais: configuração de mão (CM), ponto de articulação (PA), movimento
(M), orientação (O) e expressões não-manuais (ENM), estas últimas relacionadas a
movimentos da face, dos olhos, da cabeça ou do tronco durante a sinalização.

Esse arcabouço teórico tem implicação direta sobre o problema técnico enfrentado
neste trabalho: como o parâmetro de movimento é constitutivo do sinal — e não
incidental —, sistemas de reconhecimento automático de Libras não podem se
limitar à classificação de poses estáticas da mão em um único \textit{frame}. É
necessário capturar a dinâmica temporal do gesto, o que justifica, do ponto de
vista linguístico, a escolha de uma arquitetura de rede neural capaz de modelar
sequências (Seções~\ref{subsec:cnn} e~\ref{subsec:lstm}).

Adicionalmente, as expressões não-manuais cumprem funções gramaticais na Libras —
marcação de referência específica, referência pronominal, partícula negativa,
entre outras (QUADROS; KARNOPP, 2004) —, o que representa uma limitação
reconhecida do escopo atual deste trabalho, concentrado no rastreamento dos
\textit{landmarks} das mãos e ainda não incorporando expressão facial (ver
Seção~\ref{subsec:limitacoes}).

\subsection{Reconhecimento de língua de sinais}

O reconhecimento automático de língua de sinais (SLR) é uma tarefa da área de
interação humano-computador que consiste em traduzir, computacionalmente, sinais
executados por um sinalizante em texto ou fala correspondente. Trata-se de um
problema mais complexo do que o reconhecimento genérico de gestos, pois exige
discriminação fina entre configurações de mão, orientações e expressões faciais
sutis, todas fortemente acopladas a dependências temporais entre os
\textit{frames} de vídeo (REGO et al., 2025).

A literatura da área costuma dividir as abordagens de captura em dois grandes
grupos: sistemas baseados em sensores vestíveis (luvas instrumentadas, sensores
de eletromiografia) e sistemas baseados em imagem/vídeo, com ou sem sensores de
profundidade (câmeras RGB, RGB-D, Kinect). Sistemas baseados em imagem, como o
desenvolvido neste trabalho, têm a vantagem de não exigir \textit{hardware}
especializado do usuário, ao custo de maior sensibilidade a variações de
iluminação, oclusão e ângulo de câmera.

\subsection{Redes Neurais Convolucionais (CNN)}
\label{subsec:cnn}

Redes Neurais Convolucionais (\textit{Convolutional Neural Networks} — CNN) são
arquiteturas de aprendizado profundo especializadas na extração de
características espaciais a partir de dados estruturados em grade, como imagens
ou, no caso deste trabalho, vetores de coordenadas de \textit{landmarks} por
\textit{frame}. Camadas convolucionais aplicam filtros deslizantes que aprendem a
detectar padrões locais — bordas, formas, relações geométricas entre pontos —
sem que essas características precisem ser definidas manualmente.

No \textit{pipeline} deste trabalho, camadas convolucionais unidimensionais
(\texttt{Conv1D}) são aplicadas de forma distribuída no tempo
(\texttt{TimeDistributed}) sobre cada \textit{frame} da sequência de
\textit{landmarks}, extraindo, para cada instante, um vetor de características que
resume a configuração espacial da mão naquele \textit{frame}, antes de a
informação ser repassada às camadas recorrentes responsáveis por modelar a
dinâmica temporal (Seção~\ref{subsec:arquitetura}).

\subsection{Redes LSTM (\textit{Long Short-Term Memory})}
\label{subsec:lstm}

Redes recorrentes do tipo LSTM (\textit{Long Short-Term Memory}), propostas
originalmente por Hochreiter e Schmidhuber (1997), foram desenvolvidas para
superar a dificuldade das redes recorrentes tradicionais em aprender dependências
de longo prazo em sequências, por meio de um mecanismo de portões
(\textit{gates}) que regula o fluxo de informação retida ao longo do tempo.

Pigou et al.\ (2018) demonstraram, em uma comparação sistemática de arquiteturas
para reconhecimento de gestos em vídeo, que modelos recorrentes como LSTM e GRU
superam consistentemente abordagens baseadas em \textit{pooling} temporal
simples, por preservarem de forma mais fiel a ordem e a duração relativa dos
movimentos — um resultado que fundamenta diretamente a escolha arquitetural
adotada neste trabalho, descrita na Seção~\ref{subsec:arquitetura}.

\subsection{MediaPipe Hands}

MediaPipe Hands é uma solução de rastreamento de mãos em tempo real desenvolvida
pelo Google, capaz de predizer, a partir de uma única câmera RGB, um esqueleto de
21 pontos-chave (\textit{landmarks}) por mão, com coordenadas $(x, y, z)$,
operando em velocidade compatível com aplicações interativas mesmo em
\textit{hardware} de consumo (ZHANG et al., 2020). A solução é composta por dois
modelos em cascata: um detector de palma, responsável por localizar a região da
mão na imagem, e um modelo de \textit{landmarks}, que opera sobre essa região
recortada para predizer os 21 pontos com precisão.

A escolha do MediaPipe como etapa de extração de características, em vez do
processamento direto do vídeo RGB bruto por uma rede convolucional 3D, segue uma
tendência observada em trabalhos recentes de reconhecimento de Libras (MOREIRA,
2026), que reduz o custo computacional do \textit{pipeline} e viabiliza a
inferência em tempo real sem depender de GPUs de alto desempenho — requisito
importante para um sistema pensado para \textit{hardware} comum, como descrito
nos objetivos deste trabalho.
